\songtitle{A Crown With Bounds}

\tag*{2026}\tag{folk-ballad}\tag{original-source}\tag{kingdom-of-columbia}
\begin{notes}
From the \emph{Kingdom of Columbia} alternate history by Carlos Thompson (2021). Lyrics by Suno based on a paragraph of a historical summary by Carlos Thompson -- the first two constitutional kings, Charles and John, and the emancipation act of 1880. Retrieved from rataflechera's Suno page (V6, 28 October 2026).
\end{notes}
\begin{style}
Dixie county folk ballad with an epic historical sweep, steady loping acoustic-guitar and banjo groove, muted handclaps and brushed snare, fiddle answering the vocal lines; verses delivered as plainspoken narrative with a weathered warm lead, then a stately choral fanfare in the chorus with layered mixed voices, low brass-like harmonies and a rising fiddle countermelody. Build from intimate porch-song storytelling into a communal, regal refrain, then strip back for the final verse and end with a held choir cadence. Organic roots recording, clear vocals, room ambience, restrained cinematic depth.
\end{style}
\begin{lyrics}
\songpart{Verse 1}
Charles crossed the water with a borrowed name,\\
Bourbon blood and a foreign claim.\\
He laid Spain’s succession down by choice,\\
Took Columbian Baptism, raised a different voice.\\
A crown was waiting, but the charter held the reins;\\
He learned the courts of Europe, far from home and rain.

\songpart{Chorus}
Raise the banner, sound the hall,\\
A crown with bounds, a law for all.\\
From Charles to John, let the record ring:\\
The people’s land, the people’s king.

\songpart{Verse 2}
Charles had sons: Charles, John, Ferdinand,\\
Three roads opening across the land.\\
John took the faith and made Columbia home,\\
While Europe’s polished chambers saw his father roam.\\
A regent from ’46, steady through the years,\\
Till 1855 brought the kingdom’s turning here.

\songpart{Chorus}
Raise the banner, sound the hall,\\
A crown with bounds, a law for all.\\
From Charles to John, let the record ring:\\
The people’s land, the people’s king.

\songpart{Verse 3}
When King Charles died, they named John king,\\
And he stayed where the bells and the river sing.\\
His power was limited, yet he used it with care,\\
A measured hand at the council chair.\\
He chose his moments, spoke when it mattered most,\\
A quiet kind of courage, not a boast.

\songpart{Bridge}
From the long debate of ’72, he carried one cause on,\\
Through patient years and stubborn doors, till the old rule was gone.

\songpart{Verse 4}
In 1880, the emancipation act was signed,\\
A promise he had championed, brought at last to time.\\
Not every chain breaks with a royal hand,\\
But a voice can help turn the law of the land.\\
So remember Charles abroad, and John who chose to stay,\\
Two first monarchs, and the bounds that shaped their day.

\songpart{Final Chorus}
Raise the banner, sound the hall,\\
A crown with bounds, a law for all.\\
From Charles to John, let the record ring:\\
The people’s land, the people’s king.\\
Raise the banner, sound it clear,\\
Columbia’s story carried year to year.

\songpart{Outro}
A crown with bounds, a law for all;\\
Columbia answers when the people call.
\end{lyrics}

